\section{Инженерное применение и границы}
Метод готовит правило реагирования для процессно-ориентированной информационной системы. Карточка фиксирует требования к наблюдениям, действие, порог, версию модели, диапазон калибровки, стрессовые сценарии, эффекты, стоимости и результаты сравнений. Владелец процесса утверждает действия и требования к обслуживанию; операционный контроллер регистрирует срабатывания, доступные ресурсы и последствия. Изменения маршрутов, персонала или состава задач, устойчивые ошибки прогноза либо выход эффекта за утверждённый диапазон требуют повторной калибровки. Эти требования дополняют проверенную вычислительную процедуру.

Допуск политики зависит от целей организации. В P0 при обеих обратных связях снижение стоимости WIP сочетается с ростом частоты затяжной перегрузки. Политики, нарушающие обязательные требования к обслуживанию, следует исключать до оптимизации стоимости. Многокритериальное предписание~\cite{multikpi} рассматривает такие компромиссы. Обоснование решения должно включать частоту тяжёлых событий и, при общих ресурсах, показатели каждого процесса наряду с общей стоимостью.

Оценка охватывает сконструированные синтетические структуры, заданные стрессы и эффекты действий. Равные веса процессов не отражают распространённость механизмов в организациях. Ресурсные соответствия и мощности задаются инженером, а не обнаруживаются автоматически. Корректная структура, покрытие поведения и обоснованные эффекты — отдельные условия доверия к решению. Исследования реальных журналов характеризуют предиктивный нагрузочный сигнал, а не пользу вмешательств.

Экстраполяция остаётся допущением. Калибровка P0 с нагрузочными возвратами достигает 19 незавершённых экземпляров, ниже порога реагирования 20 и критерия тяжёлого события 100. Такие журналы не различают механизмы, совпадающие в наблюдаемом диапазоне и расходящиеся при стрессе. Усложнение структуры не решает проблему идентификации. Если решение зависит от плохо наблюдавшейся нагрузки или неопределённого эффекта, нужны дополнительные измерения либо контролируемый пилот.

\section{Заключение}
Подготовленная реакция чувствительна к выбору поведенческих компонентов и структуры процесса. Сохранение нагрузочных зависимостей улучшает последовательный контрольный случай и среднее сравнение при нагрузочных возвратах и обеих связях с заданной структурой. Расширение также выявляет процессы с ухудшением решений и ошибки discovery, устраняющие общий выигрыш. Польза зависит от процесса и предполагаемых эффектов действий. Поэтому исполнимую модель следует оценивать по независимым исходам подготовленной реакции вместе с ошибкой прогноза и соответствием структуры.

\begin{credits}
\subsubsection{Доступность материалов и воспроизводимость.}
Дополнительный исследовательский архив содержит протоколы, код, контрольные суммы входов, численные параметры моделей, траектории, результаты сравнений и скрипты воспроизведения анализа и таблиц. В нём различаются основной эксперимент и последующие диагностические проверки; вспомогательные исследования приведены полностью. Открытые журналы событий указаны по DOI версии 1.
\subsubsection{Использование генеративного ИИ.}
Генеративный ИИ использовался при написании кода симуляторов, анализа и графиков, организации экспериментов, проверке литературы, подготовке текста, переводе и самопроверке по артефактам. Синтетические данные получены документированными симуляторами; численные утверждения основаны на выполненных опытах и сохранённых результатах, проверенных скриптами анализа. Генеративный ИИ не входит в исполняемые модели процессов. Авторы отвечают за статью и исследовательские материалы.
\subsubsection{\discintname}
Авторы заявляют об отсутствии конфликтов интересов, относящихся к содержанию этой статьи.
\end{credits}
